\documentclass[11pt,reqno]{amsart}
\usepackage[utf8]{inputenc}
\usepackage{braket,amsmath,amsfonts,amsthm,url,graphicx,url,xcolor,hyperref,subcaption, multicol,dsfont, comment}



\oddsidemargin-0cm
\evensidemargin-0cm
\textwidth16.1cm
%\topmargin1cm
%\footskip1.3cm %1.5cm
%\footheight3cm
%\textheight21.5cm %23cm
%\parindent0em
\renewcommand{\baselinestretch}{1.1}

\setlength{\parindent}{4ex}

\newtheorem{lemma}{Lemma}[section]
\newtheorem{prop}[lemma]{Proposition}
\newtheorem{theorem}[lemma]{Theorem}
\newtheorem{thmx}{Theorem}
\newtheorem{cor}[lemma]{Corollary}
\newtheorem{con}[lemma]{Conjecture}
\newtheorem{prob}[lemma]{Problem}
\newtheorem{rem}[lemma]{Remark}
\newtheorem{definition}[lemma]{Definition}
\newtheorem{frem}[lemma]{Final remark}
\newtheorem{fact}[lemma]{Fact}
\newtheorem{claim}[lemma]{Claim}
\newtheorem{assump}[lemma]{Assumption}
\newtheorem{exam}[lemma]{Example}
\newtheorem{cexam}[lemma]{Counterexample}
\newtheorem{quest}[lemma]{Question}


\newcommand{\expec}{\mathbb{E}}
\newcommand{\Hil}{\mathcal{H}}
\newcommand{\ketbra}[1]{|{#1}\rangle\langle{#1}|}


\newcommand{\Haar}{H}

\newcommand{\NN}{\mathbb{N}}
\newcommand{\MM}{\mathbb{M}}
\newcommand{\RR}{\mathbb{R}}
\newcommand{\EE}{\mathbb{E}}
\newcommand{\BB}{\mathbb{B}}
\newcommand{\CC}{\mathbb{C}}
\newcommand{\E}{\mathcal{E}}
\renewcommand{\S}{\mathcal{S}}
\newcommand{\sep}{\square}
\newcommand{\id}{\hat{1}}
\newcommand{\Id}{\mathrm{Id}}

\newcommand{\cV}{\mathcal{V}}
\newcommand{\cW}{\mathcal{W}}
\newcommand{\one}{\mathds{1}}

\newcommand{\pl}{\hspace{.1cm}}

\renewcommand{\L}{\mathcal{L}}
\newcommand{\M}{\mathcal{M}}
\newcommand{\N}{\mathcal{N}}
\newcommand{\R}{\mathcal{R}}
\newcommand{\G}{\mathcal{G}}
\newcommand{\D}{\mathcal{D}}
\newcommand{\idd}{\text{id}}
\newcommand{\Rot}{\text{Rot}}
\newcommand{\Ha}{H}
\newcommand{\tr}{\text{tr}}
\newcommand{\epow}[2]{F^{({#2})}_{{#1}}}
\newcommand{\A}{\mathcal{A}}
\newcommand{\B}{\mathcal{B}}
\newcommand{\C}{\mathcal{C}}
\newcommand{\cT}{\mathcal{T}}

\newcommand{\Dmax}{D_{\infty}}
\newcommand{\BP}{\text{BP}}
\newcommand{\polylog}{\text{polylog}}

\newcommand{\newterm}[1]{{\textit{#1}}}

\newcommand{\precsucc}[2]{\text{ } {}^{\prec {#2}}_{\succ {#1}} \text{ } }

\newcommand{\compl}[2]{\C({#1} \subseteq {#2})}

\renewcommand{\deg}[1]{\mathrm{deg}({#1})}

\DeclareUnicodeCharacter{2009}{\,} 

\title{Notes on Designs from Lindbladians}
%Shallow, Unstructured, Approximate k-Designs via Quantum Relative Entropy


%\author[N. Laracuente]{Nicholas Laracuente}
%\address{Indiana University Bloomington. Bloomington, IN. 47404} \email[Nicholas Laracuente]{nlaracu@iu.edu}

\thanks{\hspace{-5mm} Part of this work was completed as an IBM Postdoc at the University of Chicago.}



\newcommand{\mbalign}{}
\newcommand{\mbnline}{}

\newcommand{\todo}[1]{{\color{red}\textbf{TODO:} #1}}

\begin{document}

\section{Background and Previous Results}
We define the $k$-fold twirl weighted by probability measure $\mu$ on $U(d)$ to be
\begin{equation}
    \Phi_{\mu, k} := \int U^{\otimes k} \rho U^{\otimes k \dagger} d \mu(U)
\end{equation}
on an input system of dimension $d^k$.
\begin{lemma}[\cite{laracuente_quantum_2025}, Lemma 3.1] \label{lem:chainmin1}
Let $(\Phi_j)_{j=1}^m$ be a given family of quantum channels with respective $(\lambda_j)_{j=1}^m$-SDPI and decoherence-free subspace projections $(\E_j)_{j=1}^m$, $\E$ a joint decoherence-free subspace projection commuting with each $\Phi_j$ (for which $\E \E_j = \E_j \E = \E$),
%\fl{The 'joint' here seems to refer to the fact that $\E_m\circ \E = \E$? You seem to use this when applying the chain rule in \eqref{eq:chain2} below.}
and $(\Psi_j)_{j=1}^m$ a family of channels each commuting with $\E$. Then
%\fl{What are the `rotations' that you're referring to here? Do you mean the $\E_j$?}
\begin{equation*}
    \begin{split}
         D(\rho \| \E(\rho)) - D((\Psi_m \Phi_m) \dots (\Psi_1 \Phi_1) (\rho) \| \E (''))
            \mbnline \mbalign \geq \sum_{j=1}^m \lambda_j D( (\Psi_{j-1} \Phi_{j-1}) \dots (\Psi_1 \Phi_1)(\rho) \| \E_j ('')) \pl,
    \end{split}
\end{equation*}
where when $''$ appears in the second argument to relative entropy, it is equal to the first argument.
\end{lemma}
\begin{prop}[\cite{laracuente_quasi-factorization_2022}, Proposition I.6] \label{prop:decay-merge}
Let $\{\Phi^t_j : j \in 1...J \in \NN\}$ be GNS self-adjoint quantum Markov semigroups such that $\Phi_j^t = \exp(- \L_j t)$ with fixed point conditional expectation $\E_{j} = \lim_{t \rightarrow \infty} \Phi_j^t$ admitting common invariant state $\sigma$. Let $\E_{\sigma *}$ be the weighted intersection fixed point conditional expectation. Let $\Phi^t$ be the semigroup generated by $\L = \sum_j \alpha_j \L_j + \L_0$, where $\L_0$ generates $\Phi_0^t$ such that $\Phi_0^t \E_{\sigma *} = \E_{\sigma *} \Phi_0^t = \E_{\sigma *}$. If the conditional expectations $(\alpha_j)$-quasi-factorize, and each $\Phi^t_j$ has $\lambda$-(C)MLSI, then $\Phi^t$ has $\lambda$-(C)MLSI.
\end{prop}
\begin{lemma}[Chain Rule] \label{lem:chainexp}
Let $\omega$ be a density and $\E$ be a conditional expectation such that $\E(\omega) = \omega$. Then for any density $\rho$,
\begin{equation*}
D(\rho \| \omega) = D(\rho \| \E(\rho)) + D(\E(\rho) \| \omega) \pl.
\end{equation*}
\end{lemma}
Lemma \ref{lem:chainexp} is a well-known identity in quantum information. An important consequence of Lemma \ref{lem:chainexp} is that for any channel $\Phi$ and conditional expectation $\E'$,
\begin{equation} \label{eq:chain-iter}
    D(\rho \|\Psi(\rho)) + D(\rho \| \E'(\rho))
        \geq D(\E'(\rho) \| \E' \circ \Psi(\rho)) + D(\rho \| \E'(\rho))
    = D(\rho \| \E' \circ \Psi(\rho)) \pl.
\end{equation}
\begin{lemma}[\cite{laracuente_quasi-factorization_2022}, Corollary II.15. \cite{gao_complete_2025}, Lemma 2.3]  \label{lem:revconv}
Let $\rho$ be a density and $\E,\Psi$ be quantum channels such that
\begin{equation}
    (1 - \epsilon) \E \prec \Psi(\rho) \prec (1 + \delta) \E
\end{equation}
%\begin{equation}
%    (1-\epsilon) \E(\rho) + \epsilon \Phi(\rho) \leq (1 + \delta) \E(\rho)
%\end{equation}
for constants $\epsilon, \delta \in (0,1)$. Assume $\rho \in \text{supp}(\E(\rho))$. Furthermore, assume $\Psi \E = \E$. Then
\begin{equation}
D(\rho \| \Psi(\rho)) \geq \beta_{\epsilon, \delta} D(\rho \| \E(\rho))
\end{equation}
with
\begin{equation} \label{eq:simpleG} 
    \beta_{\epsilon, \delta} \geq \frac{1}{{(1 + \epsilon)(1 + \delta)}} \Big ( 1 - \frac{2 (1 + \epsilon) \delta^2}{(\epsilon + \delta)(\ln (1 + \delta / \epsilon) - 1) + \epsilon} - 4 \epsilon - \epsilon^2 \Big ) \pl.
\end{equation}
If $\epsilon = \delta$, then
\begin{equation}
    \beta_{\epsilon, \epsilon} \geq \frac{1-\epsilon}{1 + \epsilon} - \frac{\epsilon}{(1 - \epsilon) (2 \ln 2 - 1)} \pl, \text{ and } \beta_{\epsilon, \epsilon} \geq 1 - 12 \epsilon \pl.
\end{equation}
If $\epsilon \leq 1/10$, then $\beta_{\epsilon, \epsilon} \geq 1/2$.
\end{lemma}
\begin{definition} \label{def:quasifactorization}
    Let $\{\E_j\}_{j=1}^m$ be a set of conditional expectations with the common invariant state $\sigma$ and $\sigma$-weighted intersection subalgebra projection $\E$. We say that the set $(\alpha_j)$-quasi-factorizes (in other texts, has $\alpha$-approximate tensorization) if
    \begin{equation}
        \sum_j \alpha_j D(\cdot \| \E_j(\cdot)) \geq D(\cdot \| \E(\cdot))
    \end{equation}
    and that it $\alpha$-quasi-factorizes if it suffices to set each $\alpha_j = \alpha$.
\end{definition}
%\begin{assump}[Entropic Gluing] \label{assump:glue} \normalfont
%    Let $\E_{A}$ and $\E_B$ be subspace projections and $\E$ project to their intersection subspace. For a situation-dependent $\epsilon > 0$,
%    \begin{itemize}
%        \item (Immediate) $D(\rho \| \E_{A} \E_B (\rho)) \geq (1-\epsilon) D(\rho \| \E(\rho))$;
%        \item ($r$-Repeated) $D(\rho \| (\E_{A}\E_B)^r (\rho)) \geq (1-\epsilon) D(\rho \| \E(\rho))$ for $\rho$-independent $r \in \NN$.
%    \end{itemize}   
%\end{assump}
\begin{rem} \normalfont \label{rem:glue}
    A relative error analog of the gluing assumption was proven to hold for $k$-fold unitary twirls in \cite[Lemma 8]{schuster_random_2025} and a weaker version shown concurrently in \cite{laracuente_approximate_2026}. In that setting of $k$-fold random unitaries, $\E_{A} \circ \E_B \precsucc{5 k^2 / q^m}{} \E$, where $q$ is the local qudit dimension, and $A$ and $B$ are subsystems that overlap on $m$ elementary qudits ($k \times m$ taken $k$-fold). Combined with Lemma \ref{lem:revconv}, the relative error gluing Lemma implies that $D(\rho \| \E_A \E_B (\rho)) \geq \beta D(\rho \| \E(\rho)) $ with $\beta \geq 1 - 60 k^2 / q^m$. Hence with Equation \eqref{eq:chain-iter}, quasi-factorization holds with every $\alpha_j \leq \beta^{-1}$.
\end{rem}

Recall the tightened Fannes-Alicki inequality on relative entropy with respect to a subspace projection \cite[Lemma 7]{winter_tight_2016}:
\begin{lemma}[] \label{lem:wsb}
    if $\rho$ and $\omega$ are states, $\epsilon := \frac{1}{2} \| (\Phi \otimes \Id - \Psi \otimes \Id)(\rho) \|_1$, and $\E$ projects to a convex subspace, then
    \begin{equation} \label{eq:wsb}
        |D(\rho \| \E(\rho)) - D(\omega \| \E(\omega))|
            \leq \epsilon \sup_{\sigma} D(\sigma \| \E(\sigma)) + (1+\epsilon) h \Big ( \frac{\epsilon}{1+\epsilon} \Big ) \pl.
    \end{equation}
\end{lemma}

\section{General: Tree Decomposition}
Consider a Lindbladian $\L = \sum_{j=1}^m \L_j$. The overarching goal of this section is to derive as general as possible a way to obtain a decay constant for $\L$ in terms of (1) the decay constants of its constituents (2) how closely compositions of respective fixed point projections approximate their intersection subspaces. In particular, by recursively applying quasi-factorization, one may often obtain efficient decay estimates for the whole in terms of sums of their parts.

We decompose $\L$ via an ordered tree $T$, where the root of $T$ holds $\Phi$ or $\L$, and each parent node holds the composition of channels or Lindbladians included in its children. The leaves of the tree correspond to smaller sub-Lindbladians for which we need a prior decay constant estimate. %  for which we have some a priori estimate of decay constants. The tree can be redundant - the children of a node need not hold disjoint subsets or subsequences. For Lindbladians, redundancy is almost trivial via the identity $\L_j = \sum_r \kappa_r \L_j / (\sum_r \kappa_r)$ for any positive family $(\kappa_r)$. For sequential channels, overlaps must be restricted to neighboring siblings such that sequential order is preserved up to the overlapping portion. The ability to insert redundancy is used in \cite[Unpublished Revision]{laracuente_quantum_2025} to increase the overlaps between fixed point conditional expectations. We assume compatible entropy decay in that each $\Phi_j$ or $\L_j$ has a non-zero decay constant, and any subsequence (resp subset sum) has a well-defined decoherence-free subspace projection.
By $\L[x]$ we respectively denote the sub-Lindbladian corresponding to node $x$, by $\E[x]$ its conditional expectation, and by $\lambda[x]$ its decay constant. By $\E[x_1, \dots, x_\ell]$ we denote the intersection subspace projection of all nodes listed (assuming it is indeed well-defined). The key is the following observation:
\begin{lemma} \normalfont \label{lem:lindblad-tree}
    For each node $x \in T$, let $y_1, \dots, y_m$ be its children.
    %For each index $i$ in the original subsequence, let $(p_{i,j} : j = 1 \dots \ell)$ be any probability distribution supported on all children in which it appears.
    If the family$(\E[y_j])_{j=1}^m$ has $(\alpha_j)_{j=1}^{m}$-quasi-factorization, and each $\L[y_j]$ has $\lambda \times \alpha_j$-CMLSI, then $\L[x]$ has $\lambda$-CMLSI.
    %If the immediate case of Assumption \ref{assump:glue} holds for each $\E[y_{j+1}]$ and $\E[y_1, \dots, y_j]$ with contant $\epsilon_j$ for each $j \in 1, \dots, \ell-1$, then $\Phi[x]$ or $\L[x]$ has decay constant at least
    %\begin{equation}
    %    \left ( \prod_{j=1}^{\ell-1} (1 - \epsilon_j) \right )
    %        \times \left ( \min_{i,j} \lambda[y_j] p_{i,j} \right ) \pl.
    %\end{equation}
\end{lemma}
Lemma \ref{lem:linblad-tree} follows almost immediately from \ref{prop:decay-merge}. To use Lemma \ref{lem:lindblad-tree} requires two main criteria:
\begin{itemize}
    \item Quasi-factorization holds between different sub-Lindbladians' fixed point projections.
    \item For the leaf subsequences, there is an a priori decay rate estimate.
\end{itemize}
We may think of Lemma \ref{lem:tree} as reducing the size of a problem efficiently. The cost of converting from spectral gaps on the leaf subsequences and then building up relative entropy decay through the tree is often much smaller than that of converting spectral gaps globally.
\begin{prop} \label{prop:tree-whole}
    Let $\L = \sum_{j=1}^m \L_m$ be a Lindbladian with tree decomposition $T$. Assume each $\L_j$ has decay constant at least $\lambda$.
    %and Assumption \ref{assump:glue} holds with constant $\epsilon_{\ell}$
    For each non-leaf node $x$ at level $\ell$, assume for the children $y_1, \dots, y_{r}$ that $\E[y_1], \dots, \E[y_r]$ $\alpha_\ell$-quasi-factorize. % $\E[y_1, \dots, y_r]$ and $\E[y_{r+1}]$ when $y_1, \dots, y_{r+1}$ are siblings at level $\ell$.
    Then $\L$ has decay constant at least $\lambda / \prod_{\ell=1}^{h_T} \alpha_{\ell}$, where $h_T$ is the height of the tree.
\end{prop}
\begin{proof}
    Assume that the $\ell$th level of $T$ contains nodes $x_1, \dots, x_{m[\ell]}$. By Lemma \ref{lem:lindblad-tree}, the Lindbladian corresponding to each node in the next level up will have minimum decay constant at least $\min_{j = 1 \dots m[\ell]} \lambda[x_j] / \alpha_\ell$. Applying iteratively completes the Proposition.
\end{proof}
\begin{theorem} \label{thm:regular-tree}
    Let $\L = \sum_{j=1}^m \L_m$ be a Lindbladian with tree decomposition $T$. Assume quasi-factorization with constant at most $\exp(c / q^r)$ when $r \geq 2$ nodes are shared between adjacent siblings, and at the $\ell$th level starting from the bottom, at least $\log_q (c) + \log_q a + \ell - 1$ elementary units are shared for some constant $a \in \RR^+$.
    Assume each $\L_j$ has decay constant at least $\lambda$. Then $\L$ has decay constant at least $\exp(- 1 / a(q-1)) \lambda$.
\end{theorem}
\begin{proof}
    By Proposition \ref{prop:tree-whole} and the summation of geometric series, if $T$ has height $h$, then $\L$ has decay constant at least
    \begin{equation}
        \lambda \times \exp \bigg ( - \sum_{\ell=1}^h c / q^{ \log_q c + a + \ell } \bigg )
            \leq \lambda \times \exp \bigg ( - \frac{1}{a} \frac{1}{q-1} \bigg ) \pl,
    \end{equation}
    which is the desired bound.
\end{proof}

\subsection{Lindbladians for Random Circuits}
Consider a system of $n$ qubits and associated connectivity hypergraph $G$. Let $|G|$ denote the number of edges in $G$. We start with the model
\begin{equation} \label{eq:circuit-hypergraph-lind}
    \L^{(k)}_G := \sum_{e \in G} \frac{|e|}{d_e} \L_e^{(k)},
    \pl \L_e^{(k)}(\rho) := \rho - \Phi_e(\rho) ,
    \pl \Phi_e(\rho) := \int U^{\otimes k} \rho U^{\dagger \otimes k} d \nu_e(U) \pl,
\end{equation}
where each $d \nu_e$ denotes a \textit{local} measure on unitaries acting on the qudits indexed by edge $e$ (acting as identity on the rest of the system), $|e|$ is the number of nodes joined by $e$, and $d_e$ is the sum of degrees of nodes connected by $e$. We will assume that each $\nu_e$ induces $\lambda_e$-decay toward the local $k$-fold twirl on its constituent qudits. Since each $\L_e^{(k)}$ generates state replacement, convexity of relative entropy implies that it has CMLSI constant 1. If $G$ is a graph, then we may think of each $U_e$ as acting on a 2-qudit subsystem. For example, the Lindbladian analog of 1-D brickwork on $n$ qudits of unspecified local dimension would have $G = \{(1,2), (2,3), \dots, (n-1, n)\}$. However, we are primarily interested in more connected hypergraphs that support $\Omega(1)$ decay rates in $n$.

The ``random circuit Lindbladian'' $\L^{(k)}_G$ is named because at a given time $t$,
\begin{equation}
    e^{- t \L}(\rho) = \sum_{\ell=0}^\infty \frac{e^{-t|G|} t^{\ell}}{\ell!} \sum_{e_1, \dots, e_\ell \in G} \Phi_{e_\ell} \circ \dots \circ \Phi_{e_1} \pl.
\end{equation}
Each sequence $\Phi_{e_\ell} \circ \dots \circ \Phi_{e_1}$ can be thought of as a random circuit that applies $\ell$ gates, each given by a randomly chosen sequence of hypergraph edges. The Lindbladian effectively applies random circuits with Poisson-distributed length. Since $|G|$, the number of edges in $G$, multiplies the probability of an event in each time interval, factors of $|G|^{\ell}$ in the numerator and denominator cancel. \todo{Check that the distribution is truly Poisson, currently this is merely assumed.} \todo{Question for Nima: is this Lindbladian equivalent to the ``Brownian circuit'' model? The answer might give some guidance on how to convert Brownian SYK to a similar form.}

First, we show that the usual ``base estimate'' $O(n)$ time for relative error convergence holds for qubit hypergraph quantum circuits as well as graphs. Subsequent steps will convert this estimate into an $n$-independent rate via recursive merging.
\begin{lemma} \label{lem:hypergraph-to-graph}
    Let $\L$ be as in Equation \eqref{eq:circuit-hypergraph-lind} with hypergraph $G$. Let $G'$ be a graph constructed by adding for each edge in $G$ an arbitrary, connected subgraph in $G'$ with at most $m_e$ edges. Then the quasi-factorization constant corresponding to $G$ is at most that of $G'$ times $m_e$.
\end{lemma}
\begin{proof}
    Let $e \in G$ be any edge, and let $G_e'$ denote the subgraph of $G'$ corresponding to that edge. Recall Lemma \ref{lem:chainexp}, and note that
    \begin{equation}
        D(\rho \| \E_{G'_e}(\rho)) = D(\rho \| \E_f(\rho)) + D(\E_f(\rho) \| \E_{G'_e}(\rho))
            \geq D(\rho \| \E_f(\rho))
    \end{equation}
    for any $f \in G'_e$, where $\E_{G'_e}$ denotes the twirl over qudits in $G'_e$, and each $G_f$ denotes the twirl corresponding to a single edge. Therefore
    \begin{equation}
        m_e D(\rho \| \E_{G'_e}(\rho)) \geq \sum_{f \in G_e'} D(\rho \| \E_f(\rho)) \pl,
    \end{equation}
    finishing the Lemma.
\end{proof}
\begin{lemma} \label{lem:hypergraph-base}
    Let $\L$ be a Lindbladian as in Equation \eqref{eq:circuit-hypergraph-lind} with hypergraph $G$ on $n$ qubits having maximum edge degree $d_E = O(1)$ in $n$ and $k$. Assume that $\lambda_e \geq \lambda_0$ for every $e \in G$. Further assume a random circuit of depth at most $\ell$ on $n$ qubits is constructed by applying a random unitary to the qubits indexed by each edge in $G$, where edges that share any nodes are always applied at different layers. Then $\L$ has $\lambda$-CMLSI with
    \begin{equation}
        \lambda \geq \lambda_0 /O(n k \mathrm{polylog}(k)) \pl.
    \end{equation}
\end{lemma}
\begin{proof}
    First, assume that $G$ is a graph. Recall \cite[Equation (51)]{belkin_approximate_2024}, so the spectral gap $s_*$ of any $\ell$-layer parallel random circuit made of 2-qubit gates is bounded as
    \begin{equation}
        s_* \leq 1 - (1 - e^{-1/2 C(q,k)})^{\ell - 1} \pl,
    \end{equation}
    where $C(q,k)$ depends on the convergence of 1-D brickwork and is calculated for several cases therein.  Via \cite{chen_incompressibility_2024}, it is known that $C(q,k)$ can be taken $O(\mathrm{polylog}(k))$.
    %If $d_E > 2$, consider the random unitary projection $\E_e$ corresponding to a single edge $e$. Let $v_1, \dots, v_{|e|}$ denote the nodes corresponding to $e$, indexing $|e| \leq d_E$ qubits, where $|e|$ is the number of qubits joined by $e$.
    %Via the Solovay-Kitaev theorem and known circuit approximation bounds, every unitary on $|e|$ qubits is $\delta$-approximated in operator norm by a sequence of $O(\mathrm{polylog}(1/\delta)) \times \exp(O(|e|))$ 2-qubit gates, independently from $n$ or $k$. By \cite[Theorem 1]{yada_non-haar_2025} ...
    When $C(q,k)$ is large enough (which we may ensure by enlarging it by a constant),
    \begin{equation}
        \log s_* \geq \frac{1}{4 C(q,k)^{\ell - 1}} \pl.
    \end{equation}
    By \cite[Lemma 3]{brandao_local_2016}, such a circuit is an $\epsilon'$-approximate relative error $k$-design when $\epsilon' \leq q^{2 k n} s_*^m$ after $m$ sequential applications. Therefore, for a given $\epsilon > 0$, it suffices to take
    \begin{equation} \label{eq:circuit-depth-base}
        m = \log_{s_*} (\epsilon q^{- 2 k n}) = \frac{\log_q \epsilon - 2 k n}{\log_q s_*} \leq 4 C(q,k)^{\ell - 1} (\log_q (1/\epsilon) + 2 k n) \pl.
    \end{equation}
   Consider the sum $\sum_{e \in G} D(\rho \| \E_e^{(k)}(\rho))$, where each $\E_e^{(k)}$ denotes the $k$-fold twirl over the associated graph edge. Repeatedly using Equation \eqref{eq:chain-iter},
   \begin{equation} \label{eq:quasifac-graph}
       m \sum_{e \in G} D(\rho \| \E_e^{(k)}(\rho)) \geq D \bigg ( \rho \bigg \| \bigg ( \prod_{e \in G} \E_e^{(k)} \bigg )^m (\rho) \bigg ) \pl ,
   \end{equation}
   where the product may be taken in an arbitrary order. Therefore, the family of conditional expectations $(\E_e^{(k)})_{e \in G}$ has $m/\beta_{\delta}$-quasi-factorization as in Lemma \ref{lem:revconv}. As long as each constituent in $\L$ has $\lambda_0$-decay, Proposition \ref{prop:decay-merge} therefore shows that $\L$ has $\lambda_0 \beta_{\delta}/m$-decay. Recall, however, that this stage of the proof assumed that $G$ was a graph.
    
   Now assume $G$ is a hypergraph with $d_E > 2$. Given an edge $e$, let $|e|$ denote the edge degree. If $v_1, \dots, v_{|e|}$ denote the nodes joined by $e$ in arbitrary order, let $f[1], \dots, f[|e|-1]$ denote edges in a graph connecting them and $\E^{(k)}_{f[1]}, \dots, \E^{(k)}_{f[|e|-1]}$ denote respective $k$-fold twirls on those adjacent qubit pairs. Using Lemma \ref{lem:hypergraph-to-graph},
   \begin{equation}
       (d_E-1) D \Big (\rho \Big  \| \E_e^{(k)}(\rho) \Big )
        \geq \sum_j D \Big (\rho \Big \| \E_{f[j]}^{(k)}(\rho) \Big ) \pl.
   \end{equation}
    Recalling Equations \eqref{eq:quasifac-graph} and \eqref{eq:circuit-depth-base}, the family of conditional expectations $(\
    E_e^{(k)})_{e \in G}$ has $d_E m / \beta_\delta$-quasi-factorization with $m$ as estimated previously. Therefore, $\L$ has $\lambda_0 \beta_{\delta}/m d_E$-CMLSI, and we may set $\delta$ to $1/n$ to complete the Lemma.
\end{proof}

\begin{prop} \label{prop:lindbladian-recur}
    Let $\L$ be a Lindbladian as in Equation \eqref{eq:circuit-hypergraph-lind} with hypergraph $G$ on $n$ qubits, assuming $\lambda_e \geq \lambda_0$ for every $e \in G$. Assume a recursive splitting of $G$ into two sub-hypergraphs at each level such that at the $\ell$th level, each joining involves at least $\ell + \Omega(\log k)$ qubits, and each of the base constituents acts on a system of size at most $O(\mathrm{polylog}(k))$. Then $\L$ has CMLSI constant at least $\lambda_0 / O(k \mathrm{polylog}(k))$ that is $\Omega(1)$ in $n$.
\end{prop}
\begin{proof}
    The Proposition will follow from Theorem \ref{thm:regular-tree} after choosing an appropriate leaf size $m$. From Remark \ref{rem:glue}, we obtain a quasi-factorization constant of $(1 - 60 k^2 / 2^{m + \ell})^{-1}$ assuming that at the lowest level, at least $m$ qubits are shared across any split. Let $m := \log_2 (240 k^2)$ so that the quasi-factorization constant is at most $(1 - 1/2^{\ell+2})^{-1} \leq \exp(- 1 / 2^{\ell+1})$. If the number of joining qubits is at first two small by a constant factor, replace the leaf level by a higher level for which the additional $\ell$ yields a sufficiently large leaf size. First lower-bound the leaf CMLSI constants via Lemma \ref{lem:hypergraph-base}, then apply Theorem \ref{thm:regular-tree}.
\end{proof}

\begin{rem} \normalfont \label{rem:convex-graph}
    The Lindbladians of Equation \eqref{eq:circuit-hypergraph-lind} can be extended to Lindbladians of the form $\L = \int \L_x d \mu(x)$, where each $x$ is a label chosen according to probability measure $\mu$, and every $\L_x$ is a Lindbladian as in Equation \eqref{eq:circuit-hypergraph-lind} with uniformly upper-bounded hypergraph degree and lower-bounded CMLSI constant.
\end{rem}

\begin{exam}[Random Circuit Lindbladian on Grid] \normalfont \label{exam:grid}
Assume that $D \geq 2$, and $n = 2^{x D}$ for some $x \in \N$. Let $G$ be a $D$-dimensional lattice hypercube graph. The lattice hypercube structure enables recursively splitting the graph in half, obtaining subgraphs of size $2^{x D - 1}$, then $2^{x D - 2}$, and so on until reaching size $m$ for some $m \in \NN$ we will set later. Successive splits will iterate through axes, so every $D$ splits, we again obtain (smaller) hypercubes. This way, we also see that at each split, the number of edges crossing the split is divided by 2 (conversely if building up, each join involves doubly many edges as its children). These edges can be assigned somewhat arbitrarily to one side of the split or the other. The full recursive splitting requires exactly $\log_2 (n/m)$ layers of recursion. The face size of an $m$-length hyper-rectangle is at least $m^{(D-1)/D}$, which lower-bounds the number of qubits at each interface. By Proposition \ref{prop:lindbladian-recur}, $\L$ has $1/O(k \mathrm{polylog}(k))$-CMLSI.

By Remark \ref{rem:convex-graph}, the example extends to convex combinations of grids such as the $n$-vertex complete graph.
\end{exam}

\begin{exam}[Complete Hypergraph] \normalfont \label{exam:complete-hyper}
    Consider the complete hypergraph $G$ of edge degree $d_E$ inducing $n$-qubit Lindbladian $\L$ as in Equation \eqref{eq:circuit-hypergraph-lind}. Via Remark \ref{rem:convex-graph} and the symmetry of this configuration, it suffices to show an equivalent estimate for an essentially arbitrary $d_E$-regular hypergraph. For any $n \in \NN$, one may naturally extend an $n \times n$ 2-D grid to a $d_E$-regular hypergraph by adding $d_E - 2$ vertices to each hypergraph edge, each of which does not directly connect to any other vertex. This produces a grid-like graph, where each would-be edge is instead replaced by a $d_E$-length nearest-neighbor line. Using a similar recursive division as in \ref{exam:grid}, the number of gates crossing each partition is similar. Therefore, we again obtain $1/O(\mathrm{polylog}(k))$-CMLSI for this extended grid, and by symmetry, for the complete hypergraph.
\end{exam}
Currently, both Examples \ref{exam:grid} and \ref{exam:complete-hyper} require some extra constraints on the number of qubits involved, since the graphs must fit certain grid symmetries. This is almost certainly an artifact of the method that can be removed with some painstaking but straightforward analysis. Whether that analysis is worth undertaking depends on whether we ultimately use these examples.



\subsubsection{Pauli Basis for Qubit Random Unitaries}

Recall the Pauli basis of qubit matrices, $(Z_j)_{j=1}^4 = (I, X, Y, Z)$. Every unitary on $n$ qubits has the form
\begin{equation}
    U = \sum_{j_1, \dots, j_n = 1}^4 c_{j_1, \dots, j_n} Z_{j_1} \otimes \dots \otimes Z_{j_n} \pl.
\end{equation}
Therefore, each $k$-fold unitary has the form
\begin{equation}
    U^{\otimes k} = \sum_{j^1_1, \dots, j^1_n = 1}^4 \dots \sum_{j^k_1, \dots, j^k_n = 1}^4 \prod_{a=1}^k c_{j^a_1, \dots, j^a_n} Z_{j^a_1} \otimes \dots \otimes Z_{j^a_n} \pl.
\end{equation}
One question is whether the integral
\begin{equation}
    \int U^{\otimes k} \rho U^{\otimes k \dagger} d\mathrm{Haar}(U)
\end{equation}
has any special Pauli form for general $k$ and $\rho$. For example, when $k=1$, it is equivalent to the sum
\begin{equation}
    \frac{1}{4^n} \sum_{j^1_1, \dots, j^1_n = 1}^4 Z_{j_1} \otimes \dots \otimes Z_{j_n} \rho Z_{j_1} \otimes \dots \otimes Z_{j_n}
\end{equation}
via the self-adjointness of the Pauli matrices and known form of qubit depolarizing noise. Some results for $k=2$ might exist in \cite[Supplemental Material]{angrisani_classically_2025}.

\section{Spin-SYK}
Consider: \url{https://arxiv.org/pdf/2509.04075}, \url{https://journals.aps.org/prb/abstract/10.1103/PhysRevB.109.094206}, \url{https://link.springer.com/article/10.1007/JHEP05(2024)280#Ack1}, \url{https://journals.aps.org/pre/abstract/10.1103/3gry-qtl3}

\section{Fermions}
\begin{con}
    The $n$-fermion Brownian SYK Lindbladian has CMLSI rate $O(1)$ in $n$.
\end{con}
The Brownian SYK model is equivalent to a Linbdladian given by
\begin{equation}
\begin{split}
     G_{eff}(\rho) = &  - \frac{1}{2} \sum_I [X_I, [X_I, \rho]] 
    = - \frac{1}{2} \sum_{i_1, \dots, i_q } \sum_{a,b}
        [\psi_{i_1}^{(a)} \dots \psi_{i_q}^{(a)}, [\psi_{i_1}^{(b)} \dots \psi_{i_q}^{(b)}, \rho] ]
    \\ = & - \frac{1}{2} \sum_{i_1, \dots, i_q } \sum_{a,b} \Big ( 
        \psi_{i_1}^{(a)} \dots \psi_{i_q}^{(a)} \psi_{i_1}^{(b)} \dots \psi_{i_q}^{(b)} \rho
            + \rho \psi_{i_1}^{(b)} \dots \psi_{i_q}^{(b)} \psi_{i_1}^{(a)} \dots \psi_{i_q}^{(a)}
    \\ & - \psi_{i_1}^{(a)} \dots \psi_{i_q}^{(a)} \rho  \psi_{i_1}^{(b)} \dots \psi_{i_q}^{(b)}
        -  \psi_{i_1}^{(b)} \dots \psi_{i_q}^{(b)} \rho \psi_{i_1}^{(a)} \dots \psi_{i_q}^{(a)}
    \Big ) \pl,
\end{split}
\end{equation}
where $a, b = 1 \dots k$ and index the copy to which each fermion operator applies. Equivalently,
\begin{equation}
\begin{split}
    G_{eff} = \sum_{i_1, \dots, i_q} \L[i_1, \dots, i_q] \pl,
\end{split}
\end{equation}
where each
\begin{equation}
\begin{split}
    \L[i_1, \dots, i_q](\rho) = & - \frac{1}{2} \sum_{a,b} \Big ( 
        \psi_{i_1}^{(a)} \dots \psi_{i_q}^{(a)} \psi_{i_1}^{(b)} \dots \psi_{i_q}^{(b)} \rho
            + \rho \psi_{i_1}^{(b)} \dots \psi_{i_q}^{(b)} \psi_{i_1}^{(a)} \dots \psi_{i_q}^{(a)}
    \\ & - \psi_{i_1}^{(a)} \dots \psi_{i_q}^{(a)} \rho  \psi_{i_1}^{(b)} \dots \psi_{i_q}^{(b)}
        -  \psi_{i_1}^{(b)} \dots \psi_{i_q}^{(b)} \rho \psi_{i_1}^{(a)} \dots \psi_{i_q}^{(a)}
    \Big ) \pl.
\end{split}
\end{equation}


\section{Local Subalgebras and Observables}
A key consequence of relative entropy contraction is that it controls diamond norm distance, an operational measure of distinguishability, through Pinsker's inequality. The usual relation that $2 \| \rho - \sigma \|_1^2 \leq D(\rho \| \sigma)$ becomes potentially weak for large $n$ and trivial as $n \rightarrow \infty$, since the initial relative entropy could be arbitrarily large. This limitation is physical - existing analyses have shown that the depth required for additive error convergence of random quantum circuits is lower bounded as $\Omega(\log n)$ \cite{dalzell_random_2022, deshpande_tight_2022}. Therefore, even $\Omega(1)$ entropy contraction cannot yield general, constant-depth additive error convergence. Instead, entropy contraction bounds the \textit{average} trace on local subalgebras via the superadditivity of relative entropy on tensor splits and Pinsker's inequality:
\begin{lemma} \label{lem:throwaway}
    Consider channel $\Phi_{\mathrm{Haar}, k}$ acting on the entirety of a $k$-fold bipartite system $A^k B^k$, followed by a partial trace of the $B$ systems. Let $\E_{dep}^A$ denote the completely depolarizing channel on $A^k$. Assuming $|B| \geq 2 k^2$, $\tr_B \circ \Phi_{\mathrm{Haar}, k} \precsucc{\epsilon}{} \E_{dep}^A$ with $\epsilon \leq 2 k^2 / |B|$.
\end{lemma}
\begin{proof}
    First consider the scenario in which $B \cong B_0 \otimes A^{n-1}$ for some $n \in \NN$ and each of the $k$ copies. By Schur-Weyl duality and using the Bell pair transpose trick to apply $\Phi_{\mathrm{Haar},k}$ to both sides,
    \begin{equation}
        D((\Id \otimes \Phi_{\mathrm{Haar},k}) (\mathrm{BP}) \| (\Id \otimes \E_{dep})('')) = O(k \log k)
    \end{equation}
    with no $n$-dependence, where $\Id$ is the identity channel acting on an auxiliary system $A^{'k} B^{'k}$ of the same type as the input, and $\E_{dep}$ is a completely depolarizing channel on system $A^k B^k$. Note that $(\Id \otimes (\E_{dep} \circ \Phi_{\mu,k}))(\mathrm{BP}) = (\Id \otimes \E_{dep})(\mathrm{BP}) = \id / |A A' B B'|^k$ is a product state on all subsystem divisions. By superadditivity of the relative entropy with respect to tensor products, letting $A_j^k$ denote the $k$-fold system corresponding to the $j$th of $n$ copies of $A^k$,
    \begin{equation}
        \sum_j D((\Id \otimes (\E_{A_j} \circ \Phi_{\mathrm{Haar},k})) (\mathrm{BP}) \| (\Id \otimes \E_{dep})('')) 
            \leq D((\Id \otimes \Phi_{\mathrm{Haar},k}) (\mathrm{BP}) \| (\Id \otimes \E_{dep})('')) \pl.
    \end{equation}
    After applying $\Phi_{\mathrm{Haar}, k}$, the resulting state is invariant under exchanging the $n$ copies of $A^k$. Hence, the individual $A_j$ terms as above are all equal to the average, $D((\Id \otimes \Phi_{\mathrm{Haar},k}) (\mathrm{BP}) \| (\Id \otimes \E_{dep})('')) / n$. To complete the first part of the Lemma, note for any $\delta > 0$ there is an $n$ sufficiently large that the aforementioned is at most $\delta$.

    Let $\Phi_{\mathrm{Haar},k}[BC]$ denote the $k$-fold twirl on $B^k C^k$ (acting as identity on $A$), where $C \cong A^n$ for an $n \in \NN$ to be determined, $\Phi_{\mathrm{Haar},k}[AB] := \Phi_{\mathrm{Haar},k} \otimes \Id^C$, and $\Phi_{\mathrm{Haar},k}[ABC]$ denote the tripartite twirl. By the SHH gluing Lemma \cite[Lemma 3]{schuster_random_2025} (see also \cite{laracuente_approximate_2026}), $\tilde{\Phi}_{\mathrm{Haar},k} \circ \Phi_{\mathrm{Haar},k}$ is a $\kappa$-approximate relative $k$-design with $\kappa \leq 2 k^2 / |B|$ as long as $|B| \geq 2 k^2$. Let $\tilde{E}_A : A^k B^k C^k \rightarrow A^k$ denote the restriction map, and note that
    \begin{equation}
        \Phi_{\mathrm{Haar},k}[BC] \circ  \Phi_{\mathrm{Haar},k}[AB]
            \precsucc{\kappa}{} \Phi_{\mathrm{Haar},k}[ABC] \pl,
    \end{equation}
    where $\Phi_{\mathrm{Haar},k}[BC]$ denotes the Haar twirl on the $BC$ system.
    %Assuming that $\kappa \leq 1/2$, $1/(1 - \kappa) = 1 + \kappa + \kappa \times \kappa/(1+\kappa) \geq 1 + 4 \kappa / 3$. Conversely, $1/(1 + \kappa) \leq 1 - \kappa$. Reversing the order inequality and applying $\tr_{BC}$ to both sides,
    %\begin{equation}
    %        \tr_{BC} \circ \Phi_{\mathrm{Haar},k}[ABC] \precsucc{4 \kappa / 3}{} \tr_{BC} \circ \Phi_{\mathrm{Haar},k}[BC]  \circ \Phi_{\mathrm{Haar},k}[AB] \pl.
    %\end{equation}
    Let $F_\omega : AB \rightarrow ABC$ denote the extension by an arbitrary state $\omega$ on $C^k$.
    \begin{equation}
    \begin{split}
        \tr_{BC} \circ \Phi_{\mathrm{Haar},k}[BC] \circ \Phi_{\mathrm{Haar},k}[AB] \circ F_\omega
            = \tr_{BC} \circ \Phi_{\mathrm{Haar},k}[BC] \circ F_\omega \circ \Phi_{\mathrm{Haar},k}
            = \tr_{B} \circ \Phi_{\mathrm{Haar},k} \pl,
    \end{split}
    \end{equation}
    since the inserted extension and middle twirl act only on systems that will subsequently be discarded. Therefore,
    \begin{equation}
        \tr_{B} \circ \Phi_{\mathrm{Haar},k} \precsucc{\kappa}{} 
            \tr_{BC} \circ \Phi_{\mathrm{Haar},k}[ABC] \circ F_\omega \pl.
    \end{equation}
    Taking $n$ arbitrarily large, $\tr_{BC} \circ \Phi_{\mathrm{Haar},k}[ABC]$ becomes an arbitrarily good approximation to $\E_{dep} \otimes \tr_{BC}$ in relative entropy. By the faithfulness of relative entropy and the Choi-Jamiolkowski isomorphism, these are therefore the same channel. Hence,
    \begin{equation}
        \tr_{B} \circ \Phi_{\mathrm{Haar},k} \precsucc{\kappa}{} 
            \E_{dep} \otimes \tr_{BC} \circ F_\omega = \Id \otimes \E_{dep} \pl,
    \end{equation}
    completing the Lemma.
\end{proof}
\todo{We could use a version of this (and underlying it, the gluing Lemma) for fermions. However, in large-$N$, it's possible that much weaker/slower convergence could hold without the gluing Lemma.}
\begin{theorem}[Average Subsystem Convergence] \label{prop:avg}
    On $n$ qudits $A_1, \dots, A_n$ of local dimension $q$, for any input state $\rho$,
    \begin{equation*}
    \begin{split}
        \frac{1}{n} \sum_{j \in 1}^n D \left (\rho^{A' A_j} \big \| \rho^{A'} \otimes \id / q \right)
        \leq & \frac{D(\rho \| (\Id \otimes \Phi_{\mathrm{Haar},k})(\rho))}{n (1 - \log_q (24 n k^2) / n)(1 - 1/n)} \pl.
    \end{split}
    \end{equation*}
    where the identity is on an auxiliary system of the same dimension denoted $A'$.
\end{theorem}
\begin{proof}
    Via data processing, relative entropy is non-increasing when discarding systems. Assume we discard a subsystem $B$ of the joint system we call $AB = A_1 \dots A_n$. Let $A'$ denote the auxiliary. By Lemma \ref{lem:throwaway},
    \begin{equation}
        \tr_B \circ \Phi_{\mathrm{Haar},k} \precsucc{\epsilon}{} \E_{dep} \circ \tr_B \pl,
    \end{equation}
    where $\E_{dep}$ is a completely depolarizing map, and $\epsilon \leq 2 k^2 / |B|$. Using \ref{lem:revconv},
    \begin{equation}
        D(\rho \| (\Id \otimes \Phi_{\mathrm{Haar},k})(\rho)) \geq 
            \beta_{\epsilon} D(\tr_{B}(\rho) \| \rho^{A'} \otimes \id/|A|) \pl.
    \end{equation}
    Assume $|A| = q^m$ and without loss of generality consider it to be the first $m$ indexed subsystems. 
    For each $j = 1 \dots m$, using \cite[Proposition 2]{capel_superadditivity_2018}
    \begin{equation}
        D(\rho^{A_j} \| \id / |A_j|) = D(\rho^{A' A_j} \| \rho^{A'} \otimes \id / |A_j|) - I(A' : A_j)_\rho \pl.
    \end{equation}
    Using \cite[Proposition 2]{capel_superadditivity_2018},
    \begin{equation}
    \begin{split}
        & D(\rho^{A'A} \| \rho^{A'} \otimes \id/|A|)
        \\ & = I(A'B' : A)_\rho + D(\tr_{B}(\rho) \| \id/|A|)
        \geq I(A' : A)_\rho + \frac{1}{m} \sum_{j = 1}^m D(\rho^{A_j} \| \id / |A_j|)
        \\ & \geq \frac{1}{m} \sum_{j=1}^m \ D(\rho^{A' A_j} \| \rho^{A'} \otimes \id / |A_j|)
            - \frac{1}{m} \sum_{j=1}^m I(A' : A_j)_\rho + I(A' : A)_\rho \pl.
    \end{split}
    \end{equation}
    The same inequality holds regardless of which subsystems were actually chosen to be the $m$ kept, so by averaging over that choice,
    \begin{equation}
    \begin{split}
        \frac{1}{n} \sum_{j=1}^n \ D(\rho^{A' A_j} \| \rho^{A'} \otimes \id / |A_j|)
            \leq & \frac{1}{m} (1 / \beta_{\epsilon}) D(\rho \| (\Id \otimes \Phi_{\mathrm{Haar},k})(\rho)) 
            \\ & + \EE_A \left [ \frac{1}{m} \sum_{j=1}^m I(A' : A_j)_\rho - I(A' : A)_\rho \right ] \pl,
    \end{split}
    \end{equation}
    where $\EE_A$ denotse the average over the choice of $A$ system, and the $j = 1 \dots m$ sum runs over indices that form $A$. Note that each $I(A : A_j)_\rho \leq I(A' : A)$, upper-bounding the average as well. To fix the choice of $m$ given that $\beta \geq 1 - 12 \epsilon \geq 1 - 24 k^2 / q^{n-m}$, let $m := n - \log_q (24 n k^2)$, so $24 k^2 / q^{n-m} \leq 1/n$.
\end{proof}
\begin{cor}
    Let $\Phi$ be a quantum channel with CMLSI constant $\lambda$ toward fixed point projection $\Phi_{\mathrm{Haar},k}$. Let $\E_1, \dots, \E_{n}$ be a family of restriction maps to disjoint subsystems each of local dimension $q$. Then
    \begin{equation*}
        \frac{1}{n} \sum_{j=1}^{n} \left \| \E_j \circ \Phi - \E_{dep} \right \|_{\Diamond}^2 \leq \frac{ 4 k \exp(-\lambda /2) \log q }{(1 - \log_q (24 n k^2) / n)(1 - 1/n)} = 4 k \exp(-\lambda / 2) \log q  + O \Big ( \frac{1}{n} \Big ) \pl,
    \end{equation*}
    where $\E_{dep}$ replaces its input by $\id / q$, the complete mixture.
\end{cor}
\begin{proof}
    Observe that by the data processing inequality, we may trace out any extra qudits, obtaining an effective system of $n$. Apply Proposition \ref{prop:avg}, bunching qudits into units of size $m$, noting that the maximum, global relative entropy is $2 n k \log q$. Apply Pinsker's inequality following , 
\end{proof}
The next step is to apply:
\begin{enumerate}
    \item Light cone bounds to the initial state.
    \item Translation invariance to the final state.
\end{enumerate}
With these, we should be able to show that \textit{each individual} local observable is close to the Haar-twirled value.
\begin{cor} \label{cor:local}
    Let $\Phi$ be a quantum channel with CMLSI constant $\lambda$ toward fixed point projection $\Phi_{\mathrm{Haar},k}$ on system $AB$. Let $\E_1, \dots, \E_{n}$ be a family of restriction maps to disjoint subsystems each of local dimension $q$ all contained in the light cone of input subsystem $A$, and let $R$ denote an arbitrary reference or auxiliary system. Assume that at the output, the ratio of the largest entropy to the average is at most $\alpha$. Then for each $j = 1 \dots n$ and any input $\rho$,
    \begin{equation}
    \begin{split}
    D \left ((\Id^R \otimes \E_j)(\rho) \big \| \rho^{R} \otimes \id / q \right)
    & \leq \frac{\alpha D(\rho^{RA} \otimes \id /|B| \| (\Id^R \otimes \Phi_{\mathrm{Haar},k})(\rho^{RA} \otimes \id / |B|))}{n (1 - \log_q (24 n k^2) / n)(1 - 1/n)}
    \\ & \leq \frac{2 \alpha \exp(-\lambda) \log |A|}{n (1 - \log_q (24 n k^2) / n)(1 - 1/n)} \pl,
    \end{split}
    \end{equation}
    and
    \begin{equation}
        \left \| \E_j \circ \Phi - \E_{dep} \right \|_{\Diamond}^2
        \leq \frac{4 \alpha \exp(-\lambda / 2) \log |A|}{n (1 - \log_q (24 n k^2) / n)(1 - 1/n)} \pl.
    \end{equation}
\end{cor}

\section{Other Considerations}
\begin{itemize}
    \item If we consider 2605.22909, would our results show that all-to-all Brownian and conventional random circuit ensembles converge to approximately the same thing, thereby confirming that nonlinear XEB suffices more generally (and also without the all-to-all assumption)?
    \item According to Max West, even a single `bit' symmetry might support the no-go argument if it leads to an observable, positive operator. The idea is to use an observable to replace the pure state. Apparently West tried this, and the bound was vacuous, so designs are not ruled out. Might relate to getting a symmetry already in the 1-design case.
\end{itemize}

\bibliographystyle{unsrt}
\bibliography{ref}

\end{document}